% CCSS 7.SP.A.1–7.SP.B.4 — Statistics: Sampling and Comparing Populations
\section{Populations and Samples}

\begin{learningGoals}
\begin{itemize}
    \item Understand the difference between a population and a sample
    \item Know why random sampling is important
\end{itemize}
\end{learningGoals}

\begin{conceptBox}{Population vs. Sample}
\begin{itemize}[leftmargin=*]
    \item \textbf{Population}: The whole group you want to study
    \item \textbf{Sample}: A smaller group chosen from the population
\end{itemize}
\end{conceptBox}

\begin{workedExample}{Choosing a Sample}
A school wants to know students' favorite lunch. They survey $50$ students out of $500$.
\answerHighlight{The $50$ surveyed students are a sample; all $500$ are the population.}
\end{workedExample}

\tipBox{A sample must be random and representative to give valid results.}

\begin{practiceBox}{Populations and Samples}
\resetProblems

\prob What is the population if a survey asks $100$ people in a city about recycling?
\answerExplain{All people in the city}{The population is everyone the survey wants to learn about.}

\prob What is a sample in a poll of $1,000$ students where $100$ are asked?
\answerExplain{$100$ students}{The sample is the $100$ students actually surveyed.}

\prob Why is random sampling important?
\answerExplain{It avoids bias}{Random sampling helps make sure the sample represents the whole population.}

\prob If a sample is not random, what problem can occur?
\answerExplain{Biased results}{Non-random samples may not represent the population, leading to incorrect conclusions.}

\prob Give an example of a population and a sample.
\answerExplain{Population: all seventh graders; sample: $50$ seventh graders}{The sample is a smaller group chosen from the population.}

\end{practiceBox}
